%============================================================
% Chapter 12: Bits in Number Theory
% STATUS: PLANNED  (generated by tools/outline.py; replace this file when the chapter is written)
%============================================================
\chapter{Bits in Number Theory}
\label{ch:number_theory}

\begin{tcolorbox}[enhanced, colback=white, colframe=secondaryblue, boxrule=1pt, arc=2mm]
\textbf{\large\textcolor{primaryblue}{Binary meets arithmetic: parity of binomial coefficients, powers of two, the 2-adic valuation and binary algorithms for gcd and powers.}}
\end{tcolorbox}

\chapterfacts{1500}{2400}{4}{8}{30}

\begin{plannedbox}[Planned Contents --- this chapter is not written yet]
\textbf{Part A --- Lessons}
\begin{enumerate}[leftmargin=*, label=\textbf{A\arabic*.}, itemsep=1pt]
    \item The 2-adic valuation $\nu_2(n)$, trailing zeros, and odd parts
    \item Lucas mod 2 and Kummer: $\binom{n}{k}$ is odd iff $k \subseteq n$; carries count the factors of 2
    \item Legendre's formula in base 2: $\nu_2(n!) = n - \operatorname{popcount}(n)$
    \item Binary exponentiation and binary multiplication (the algebra is the math notes)
    \item The binary GCD (Stein) and why it uses only shifts and subtraction
\end{enumerate}
\textbf{Part B --- Patterns} (with candidate problems, checked when the chapter is written)
\begin{enumerate}[leftmargin=*, label=\textbf{P\arabic*.}, itemsep=1pt]
    \item \textbf{Parity of binomial coefficients}
    \item \textbf{Powers of two and odd parts}
    \item \textbf{Factorials in base 2}
    \item \textbf{Binary algorithms}
\end{enumerate}
\end{plannedbox}
